\section{Méthodes numériques}
\label{sec:methodes_numeriques}

\subsection{L'approximation de grand portefeuille homogène}

La simulation Monte Carlo du portefeuille complet est coûteuse : pour chaque valeur candidate
du paramètre de corrélation, il faut générer $M$ scénarios de $N$ temps de défaut corrélés, puis
calculer les pertes de tranche pour toutes les dates de paiement. Lorsque $N$ est grand et que le
portefeuille peut être supposé homogène, une approximation analytique, due à Vasicek (2002) et
connue sous le nom de LHP (\emph{Large Homogeneous Portfolio}), permet d'éviter entièrement
la simulation.

Considérons un portefeuille de $N$ entités de même probabilité de défaut marginale $p(t) = F_i(t)$,
de même taux de recouvrement $R$, et dont la structure de dépendance est décrite par le modèle
gaussien à un facteur~\eqref{eq:facteur_gauss}. Conditionnellement au facteur systémique $M = m$,
les indicateurs de défaut $\mathbb{1}_{\{\tau_i \leq t\}}$ sont indépendants et de loi de Bernoulli
de paramètre
\begin{equation}
    q(m, t) = \Phi\!\left(\frac{\Phi^{-1}(F_i(t)) - \sqrt{\rho}\, m}{\sqrt{1-\rho}}\right).
\end{equation}
Par la loi forte des grands nombres appliquée conditionnellement à $M$, le taux de perte du
portefeuille converge presque sûrement lorsque $N \to \infty$ vers
\begin{equation}
    \ell(t) \xrightarrow[N\to\infty]{\text{p.s.}} L(t) := (1-R)\, q(M, t).
\end{equation}
La distribution des pertes dans la limite LHP est donc entièrement déterminée par celle du
facteur $M \sim \mathcal{N}(0,1)$. La perte espérée de la tranche $[K_1, K_2]$ à la date $t$ s'écrit
par le théorème de transfert
\begin{align}
    \begin{split}
    \mathrm{ETL}(t) & = \mathbb{E}^{\mathbb{P}^\ast}\!\left[\min\!\bigl(\max(L(t) - K_1, 0),\, K_2 - K_1\bigr)\right] \\
    & = \int_{-\infty}^{+\infty} \min\!\bigl(\max((1-R)q(m,t) - K_1, 0),\, K_2-K_1\bigr)\, \phi(m)\, dm,
    \label{eq:etl_lhp}
    \end{split}
\end{align}
où $\phi$ est la densité de la loi normale centrée réduite. Cette intégrale unidimensionnelle est évaluée
numériquement par une méthode de quadrature\footnote{Dans le code de simulation une méthode de quadrature
de Gauss-Legendre (voir section \ref{subsec:Gauss_Leg}) est utilisée cependant l'espérance étant ici sur une loi gaussienne une quadrature de Gauss
-Hermite pourrait être plus efficace. Néanmoins face à la régularité de la fonction la différence est négligeable
sachant que les performances sont déjà largement suffisantes.}, ce qui réduit le calcul d'une espérance sur un
portefeuille de taille arbitraire à une simple somme de quelques centaines de termes.

\subsection{Quadrature de Gauss-Legendre}\label{subsec:Gauss_Leg}

L'intégrale~\eqref{eq:etl_lhp} est calculée par une quadrature de Gauss-Legendre appliquée
après troncature du domaine d'intégration à l'intervalle $[-6, 6]$, qui capture plus de
$99.9999\,\%$ de la masse de $\mathcal{N}(0,1)$. On substitue $m \in [-6, 6]$ par le changement de
variable affine $m = 6x$ avec $x \in [-1, 1]$, puis on approche l'intégrale par
\begin{equation}
    \int_{-6}^{6} g(m)\, dm \approx 6 \sum_{k=1}^{K} w_k\, g(6 x_k),
\end{equation}
où $(x_k, w_k)_{k=1,\ldots,K}$ sont les $K$ nœuds et poids de Gauss-Legendre sur
$[-1,1]$ donnés l'algorithme de Golub-Welsh en calculant les valeurs propres et vecteurs propres 
de la matrice de Jacobi
\begin{equation}
    \mathbf{L}_K =
    \begin{bmatrix}
        0   & a_1 & 0   & \dots  & 0      \\
        a_1 & 0   & a_2 & \dots  & 0      \\
        0   & a_2 & 0   & \ddots & \vdots \\
        \vdots  & \vdots & \ddots   & \ddots & a_{n-1} \\
        0  & 0 & \dots   & a_{n-1} & 0
    \end{bmatrix},
\end{equation}
où $a_k = \frac{1}{\sqrt{4k^2-1}}$ pour $k\in[1, n-1]$. Les valeurs propres obtenues sont alors
égales aux racines du $K$-ième polynôme de Legendre $L_K(x)$ donnant les nœuds $(x_k)_{k=1,\ldots,K}$
et les poids correspondants sont donnés par
\begin{equation}
    w_k = 2\left[v^{(i)}_1\right]^2 , \quad k=1,\ldots,K,
\end{equation}
où $v^(i)_1$ désigne la première composante du $i$-ième vecteur propre normalisé de $\mathbf{L}_K$.

Nous utilisons $K = 501$ nœuds dans l'ensemble des calculs, ce qui s'avère
largement suffisant compte tenu de la régularité de l'intégrande. Les nœuds et poids sont
précalculés une seule fois en début de session et réutilisés dans toutes les évaluations, afin
de minimiser le coût des appels répétés lors de la calibration.

La quadrature de Gauss-Legendre est préférée à une grille uniforme, car à nombre de points
égal, l'erreur est exponentiellement plus petite pour des fonctions régulières. Pour le modèle
de corrélation stochastique, une double quadrature est mise en œuvre : les nœuds extérieurs
portent sur $\rho \in [0,1]$ selon la loi Bêta (31 nœuds), et les nœuds intérieurs portent sur
$m \in [-6,6]$ selon la loi normale (501 nœuds), le tout évalué simultanément sur un tenseur
de dimensions $(n_t, 31, 501)$ par opérations vectorisées.

\subsection{Inversion numérique du seuil de défaut pour la copule t-gaussienne}

Pour le modèle t-gaussien de Hull-White, le seuil de défaut $c(t)$ vérifiant
$F_{X_i}(c_i(t)) = F_i(t)$ n'est pas disponible analytiquement, la loi marginale de $X_i$ étant un
mélange de lois normales conditionnées par $\tilde{M} \sim \mathcal{T}(\nu_1)$. Le calcul de $c_i(t)$ pour
chaque date de paiement requiert donc une résolution numérique de l'équation
\begin{equation}
    G(c) := \mathbb{E}^{\mathbb{P}^\ast}_{\tilde{M}}\!\left[\Phi\!\left(\frac{c - \sqrt{\rho}\,\tilde{M}}{\sqrt{1-\rho}}\right)\right] = F_i(t),
\end{equation}
où l'espérance est approchée par quadrature de Gauss-Legendre sur la densité de Student.
Cette équation est résolue par la méthode de Newton-Raphson, dont l'itération s'écrit
\begin{equation}
    c^{(k+1)} = c^{(k)} - \frac{G(c^{(k)}) - p(t)}{G'(c^{(k)})},
    \quad G'(c) = \mathbb{E}^{\mathbb{P}^\ast}_{\tilde{M}}\!\left[\frac{1}{\sqrt{1-\rho}}\,\phi\!\left(\frac{c - \sqrt{\rho}\,\tilde{M}}{\sqrt{1-\rho}}\right)\right].
\end{equation}
L'initialisation est effectuée avec $c^{(0)} = \Phi^{-1}(F_i(t))$ (valeur gaussienne), ce qui
assure une convergence rapide en six à huit itérations environ. La procédure est vectorisée
sur toutes les dates de paiement simultanément, de sorte que le coût marginal du calcul de
$c_i(t)$ par rapport au cas gaussien demeure négligeable.

\subsection{Simulation Monte Carlo et validation}

Bien que le modèle LHP soit utilisé pour la calibration, la simulation Monte Carlo reste
indispensable pour deux usages. D'une part, elle sert de référence de validation pour évaluer
la qualité de l'approximation grand portefeuille sur le portefeuille fini de $n = 120$ noms.
D'autre part, elle est utilisée pour le modèle gaussien lors de la calibration sur la RMSE
absolue, suivant la pratique historique.

L'algorithme de simulation est identique pour les modèles gaussien et t-gaussien : on génère
$M = 10\,000$ scénarios du facteur systémique et des termes idiosyncratiques, on calcule les
temps de défaut par inversion, puis les pertes de tranche à chaque date de paiement. Les
spreads sont déduits par la condition d'absence d'arbitrage entre jambe premium et jambe
protection. Pour limiter la variabilité numérique entre deux évaluations successives de la
fonction objectif, la graine aléatoire est fixée avant chaque simulation.

Les écarts entre les spreads LHP et Monte Carlo sont inférieurs à 5 points de base pour
toutes les tranches et toutes les dates, ce qui valide l'approximation grand portefeuille pour
un portefeuille de 120 noms.

\subsection{Métriques et stratégies d'optimisation}

Pour chaque date $i$ et chaque modèle, la calibration consiste à minimiser un écart entre les
spreads de marché $(s_k^\text{mkt})_{k=1,\ldots,4}$ et les spreads du modèle $(s_k^\text{mod}(\theta))_{k=1,\ldots,4}$
sur les quatre tranches. Deux métriques ont été considérées.

\noindent La RMSE absolue, définie par
\begin{equation}
    \mathcal{L}_\text{abs}(\theta) = \sqrt{\frac{1}{4}\sum_{k=1}^{4}\bigl(s_k^\text{mod}(\theta) - s_k^\text{mkt}\bigr)^2},
\end{equation}
donne un poids dominant à la tranche equity dont les spreads sont de l'ordre de $400$ points
de base, contre moins de $2$ points de base pour la tranche super-senior. Cette métrique a été
utilisée historiquement et peut-être retrouvée dans la littérature. Elle est implémentée ici avec simulation Monte Carlo pour la copule
gaussienne, afin de reproduire les résultats de la littérature.

\noindent La RMSE logarithmique, définie par
\begin{equation}
    \mathcal{L}_\text{log}(\theta) = \sqrt{\frac{1}{4}\sum_{k=1}^{4}\ln\!\left(\frac{s_k^\text{mod}(\theta)}{s_k^\text{mkt}}\right)^2},
\end{equation}
traite toutes les tranches de manière équilibrée : une erreur multiplicative d'un facteur $r$
contribue $\ln(r)^2$ quel que soit le niveau de spread. Cette métrique est utilisée pour la
calibration LHP de la copule gaussienne et pour l'ensemble des modèles avancés.

Pour chaque modèle, la procédure d'optimisation se déroule en deux étapes. Une recherche
sur grille couvre l'espace des paramètres de manière exhaustive et fournit un point de départ
à l'optimiseur local. Pour les modèles à un paramètre (copule gaussienne LHP), la grille est
unidimensionnelle avec 200 valeurs de $\rho \in [0.01, 0.99]$. Pour les modèles à deux paramètres
(t-gaussien, corrélation stochastique), une grille rectangulaire $15 \times 15$ est utilisée.
Pour le modèle BGL à trois paramètres, une grille $10 \times 10 \times 10$ réduite à $\rho_1 \in [0.01, 0.5]$,
$\rho_2 \in [0.5, 0.99]$ et $q \in [0.5, 0.99]$ est employée, exploitant la contrainte $\rho_1 < \rho_2$.

Le raffinement local est effectué par l'algorithme de Nelder-Mead pour les modèles t-gaussien,
BGL et corrélation stochastique, du fait de sa robustesse aux discontinuités numériques légères
et de son absence de calcul de gradient. Pour la copule gaussienne LHP, le minimum de la fonction
régulière est affiné par la méthode L-BFGS-B (méthode de Broyden-Fletcher-Goldfarb-Shanno), plus rapide grâce à l'information de second ordre
approchée.

Les paramètres calibrés et les spreads modèles associés sont sauvegardés dans des fichiers
cache, ce qui permet d'éviter de relancer les calculs lors de modifications ultérieures du code
d'affichage. Le calcul des smiles de corrélation implicite (inversion numérique de la formule
LHP gaussienne sur les spreads du modèle calibré, par méthode de Brent) utilise également
ces caches.
